%% ma: L11-B09-0010
%% lop: 11
%% chuong: 3
%% bai: 9
%% dang: 11.9.8
%% loai: TL
%% muc: VD
%% dap_an: "-"
%% nguon: "de-on-tap-giua-ki-1-toan-11-nam-2026-2027-truong-thpt-ngo-quyen-da-nang.pdf (M09, Đề ôn tập giữa kì 1 Toán 11 THPT Ngô Quyền, Đà Nẵng), Mã đề 0111, Phần IV Câu 3"
%% kien_thuc: "Bài 9 (các số đặc trưng đo xu thế trung tâm)"
%% trang_thai: chinh-thuc
%% ghi_chu: "Đã duyệt 29/10/2026."
\begin{ex}
Để đánh giá mức độ chuyên cần tự học của học sinh, giáo viên chủ nhiệm ghi nhận thời gian tự học ở nhà mỗi ngày của $40$ học sinh lớp 11B trong một tuần và tổng hợp thành bảng số liệu ghép nhóm sau:
\begin{center}\begin{tabular}{|l|c|c|c|c|c|}\hline
Thời gian (giờ) & $[1;2)$ & $[2;3)$ & $[3;4)$ & $[4;5)$ & $[5;6)$\\\hline
Số học sinh & $6$ & $14$ & $12$ & $5$ & $3$\\\hline
\end{tabular}\end{center}
a) Tính thời gian tự học trung bình và trung vị của mẫu số liệu ghép nhóm trên.

b) Giáo viên đưa ra định hướng: ``Nếu thời gian tự học trung bình của lớp đạt từ $3$ giờ trở lên hoặc có từ $50\%$ số học sinh trở lên tự học từ $3$ giờ trở lên mỗi ngày thì phong trào tự học của lớp đạt mức Tốt''. Dựa vào các số liệu trên, hãy cho biết phong trào tự học của lớp 11B có đạt mức Tốt theo định hướng của giáo viên hay không?
\loigiai{a) \[\bar x=\dfrac{1{,}5\cdot6+2{,}5\cdot14+3{,}5\cdot12+4{,}5\cdot5+5{,}5\cdot3}{40}=\dfrac{125}{40}=3{,}125\text{ (giờ).}\]

$\dfrac n2=20$; tần số tích lũy: $6;20;32;37;40$, nhóm chứa trung vị là $[2;3)$: $M_e=2+\dfrac{20-6}{14}\cdot1=3$ (giờ).

b) \[\bar x=3{,}125\ge3\] (đồng thời có $12+5+3=20$ học sinh, tức $50\%$, tự học từ $3$ giờ trở lên), nên phong trào tự học của lớp 11B đạt mức Tốt.}
\end{ex}
